\documentclass[11pt,a4paper]{article}
\usepackage[utf8]{inputenc}
\usepackage[margin=1in]{geometry}
\usepackage{amsmath,amssymb,amsfonts}
\usepackage{booktabs}
\usepackage{tabularx}
\usepackage{xcolor}
\usepackage{listings}
\usepackage{hyperref}
\usepackage{fancyhdr}
\usepackage{enumitem}
\usepackage{tcolorbox}
\usepackage{titlesec}

% Colors
\definecolor{primary}{RGB}{24, 76, 120}
\definecolor{secondary}{RGB}{192, 57, 43}
\definecolor{accent}{RGB}{39, 174, 96}
\definecolor{darkgray}{RGB}{44, 62, 80}
\definecolor{lightgray}{RGB}{245, 247, 250}
\definecolor{codebg}{RGB}{240, 243, 246}

% Listings setup for SPICE
\lstdefinelanguage{SPICE}{
  morekeywords={ac, tran, step, param, meas, model, param, temp, include, subckt, ends, dec, lin, oct},
  sensitive=false,
  morecomment=[l]{*},
  morecomment=[l]{;},
  morestring=[b]",
}

\lstset{
  language=SPICE,
  backgroundcolor=\color{codebg},
  basicstyle=\ttfamily\small\color{darkgray},
  keywordstyle=\bfseries\color{primary},
  commentstyle=\itshape\color{gray},
  stringstyle=\color{secondary},
  frame=single,
  rulecolor=\color{lightgray},
  breaklines=true,
  showstringspaces=false,
  tabsize=2,
  captionpos=b
}

% Header & Footer
\pagestyle{fancy}
\fancyhf{}
\fancyhead[L]{\color{gray}\small Discrete Superheterodyne Spectrum Analyzer}
\fancyhead[R]{\color{gray}\small LTSpice Simulation \& Implementation Guide}
\fancyfoot[C]{\color{gray}\thepage}
\renewcommand{\headrulewidth}{0.4pt}
\renewcommand{\headrule}{\hbox to\headwidth{\color{lightgray}\leaders\hrule height \headrulewidth\hfill}}

% Section styling
\titleformat{\section}{\Large\bfseries\color{primary}}{\thesection}{1em}{}[\titlerule]
\titleformat{\subsection}{\large\bfseries\color{darkgray}}{\thesubsection}{1em}{}
\titleformat{\subsubsection}{\normalsize\bfseries\color{darkgray}}{\thesubsubsection}{1em}{}

% Hyperref setup
\hypersetup{
    colorlinks=true,
    linkcolor=primary,
    urlcolor=primary,
    citecolor=secondary
}

\title{
    \vspace{-1.5cm}
    {\Huge\bfseries\color{primary} Discrete Superheterodyne Spectrum Analyzer}\\[0.3cm]
    {\Large\color{darkgray} Comprehensive Circuit Implementation \& LTSpice Simulation Guide}
}
\author{\textbf{System Design \& Simulation Plan}}
\date{August 31, 2026}

\begin{document}

\maketitle
\thispagestyle{fancy}

\begin{tcolorbox}[colback=lightgray,colframe=primary,title=\textbf{Project Overview \& Scope}]
\textbf{Target Architecture:} Swept Superheterodyne Spectrum Analyzer\\
\textbf{Input Frequency Range:} 0 to 30~kHz (Audio to Low-Ultrasound)\\
\textbf{Intermediate Frequencies (IF):} 25~kHz (Channel 2) and 35~kHz (Channel 3)\\
\textbf{Resolution Bandwidth (RBW):} 200~Hz ($Q \approx 125$ to $175$)\\
\textbf{Core Hardware:} TL072 Dual JFET Op-Amps, 2N2222 Transistors, 2N3904 Transistors, 1N4148 Diodes
\end{tcolorbox}

\section{System Architecture \& Signal Chain}

The instrument follows the classic swept superheterodyne architecture. An incoming RF signal is buffered, mixed down to a fixed intermediate frequency (IF) via a swept local oscillator (LO), filtered through high-$Q$ bandpass stages, envelope detected, smoothed via a video bandwidth (VBW) filter, and displayed against the LO sweep ramp ($X$-axis):

\begin{center}
\fbox{RF Input} $\longrightarrow$ \fbox{Preamp / Attenuator} $\longrightarrow$ \fbox{2N2222 Gilbert Mixer} $\longrightarrow$ \fbox{Active IF Filter Bank} $\longrightarrow$ \fbox{Envelope Detector} $\longrightarrow$ \fbox{VBW Filter} $\longrightarrow$ \textbf{Y-OUT}\\
\vspace{0.2cm}
\hspace{4.2cm} $\uparrow$\\
\hspace{3.4cm} \fbox{Sawtooth VCO (LO)} $\longrightarrow$ \textbf{X-OUT (Tuning Ramp)}
\end{center}

\section{Block-by-Block Circuit Implementation}

\subsection{Block 1: RF Input Preamp \& Scaling Buffer}
\begin{itemize}[leftmargin=*]
    \item \textbf{Function:} Provides high input impedance ($1\text{ M}\Omega$), AC coupling, and selectable gain ($1\times, 10\times$) to scale audio signals (10~mV to 1~V$_{\text{rms}}$) to optimal mixer levels ($0.5\text{ to } 1.0\text{ V}_{\text{pp}}$).
    \item \textbf{Op-Amp:} TL072 JFET-input op-amp on $\pm 12\text{V}$ (or $\pm 9\text{V}$) rails.
    \item \textbf{Component Values:} $C_{\text{in}} = 1.0\,\mu\text{F}$, $R_{\text{in}} = 100\,\Omega$, $R_{\text{bias}} = 1.0\text{ M}\Omega$, $R_f = 9.0\text{ k}\Omega$, $R_g = 1.0\text{ k}\Omega$, $C_{\text{out}} = 1.0\,\mu\text{F}$.
\end{itemize}

\subsection{Block 2: Discrete Gilbert Cell Four-Quadrant Mixer}
\begin{itemize}[leftmargin=*]
    \item \textbf{Function:} Multiplies the RF signal with the sweeping LO to produce sum ($f_{\text{LO}} + f_{\text{RF}}$) and difference ($|f_{\text{LO}} - f_{\text{RF}}|$) products.
    \item \textbf{Topology:} 6-transistor active Gilbert cell using 2N2222 NPN BJTs with resistive bias ladder.
    \item \textbf{Biasing Ladder (from +12V rail):} $R_1 = 3.3\text{ k}\Omega$, $R_2 = 2.2\text{ k}\Omega$, $R_4 = 2.2\text{ k}\Omega$, $R_3 = 1.5\text{ k}\Omega$.
    \item \textbf{Current Tail \& Loads:} $R_{\text{tail}} = 390\,\Omega$ on $Q_{\text{Tail}}$, $R_{\text{load1}} = R_{\text{load2}} = 2.2\text{ k}\Omega$.
    \item \textbf{Decoupling \& Coupling:} $C_1, C_2, C_3 = 10\,\mu\text{F}$; LO coupling $C_4, C_5 = 100\text{ nF}$.
\end{itemize}

\subsection{Block 3: Active IF Bandpass Filter Bank}
\subsubsection{Channel 2: 25~kHz Sallen-Key Bandpass Filter}
\begin{itemize}[leftmargin=*]
    \item \textbf{Parameters:} $f_0 = 25,000\text{ Hz}$, $\text{BW} = 200\text{ Hz}$, $Q = 125$.
    \item \textbf{Components:} TL072 Op-Amps, $R_3 = R_5 = 796\text{ k}\Omega$ (1\%), $R_4 = 6.34\text{ k}\Omega$ (1\%), $C_3 = C_4 = 1.0\text{ nF}$ (C0G/NP0).
    \item \textbf{Input Attenuator:} $R_{\text{div1}} = 10\text{ k}\Omega, R_{\text{div2}} = 2.2\text{ k}\Omega$ ($-15\text{ dB}$ reduction) to eliminate passband rail saturation.
\end{itemize}

\subsubsection{Channel 3: 35~kHz Multiple Feedback (MFB) Filter (Redesign)}
\begin{itemize}[leftmargin=*]
    \item \textbf{Rationale:} The previous Sallen-Key design oscillated unconditionally due to $Q=175$ exceeding the stability limit of the op-amp. The MFB topology uses inverting feedback, ensuring unconditional stability.
    \item \textbf{Design Equations ($C_1 = C_2 = C = 1.0\text{ nF}$, $A_v = 2.0$):}
    \begin{align*}
        R_3 &= \frac{Q}{\pi f_0 C} = \frac{175}{\pi \times 35000 \times 1\times 10^{-9}} = \mathbf{1.591\text{ M}\Omega} \quad (\text{Standard } 1.58\text{ M}\Omega)\\
        R_1 &= \frac{R_3}{2 A_v} = \frac{1.591\text{ M}\Omega}{2 \times 2.0} = \mathbf{397.8\text{ k}\Omega} \quad (\text{Standard } 392\text{ k}\Omega)\\
        R_2 &\approx \frac{R_3}{4 Q^2} = \frac{1.591\text{ M}\Omega}{4 \times 175^2} = \mathbf{12.98\,\Omega} \quad (\text{Standard } 13\,\Omega \text{ or } 10\,\Omega + 3\,\Omega)
    \end{align*}
\end{itemize}

\subsection{Block 4: Diode Envelope Detector \& Buffer}
\begin{itemize}[leftmargin=*]
    \item \textbf{Components:} 1N4148 fast switching diode, $C_1 = 10\text{ nF}$, $R_1 = 22\text{ k}\Omega$ ($\tau = 220\,\mu\text{s}$).
    \item \textbf{Buffer:} TL072 configured as unity-gain voltage follower ($V_{\text{out}} = V_{\text{in}}$).
\end{itemize}

\subsection{Block 5: Video Bandwidth (VBW) Low-Pass Filter}
\begin{itemize}[leftmargin=*]
    \item \textbf{Selectable Modes:} 
    \begin{itemize}
        \item Low VBW ($f_c \approx 50\text{ Hz}$): $R = 33\text{ k}\Omega, C = 100\text{ nF}$ (Max smoothing; sweep $\ge 200\text{ ms}$).
        \item Medium VBW ($f_c \approx 200\text{ Hz}$): $R = 10\text{ k}\Omega, C = 82\text{ nF}$.
        \item High VBW ($f_c \approx 1\text{ kHz}$): $R = 10\text{ k}\Omega, C = 15\text{ nF}$ (Fast sweep response).
    \end{itemize}
\end{itemize}

\subsection{Block 6: Sawtooth Sweep Generator \& VCO}
\begin{itemize}[leftmargin=*]
    \item \textbf{Integrator:} TL072, $C_{\text{time}} = 10\text{ nF}\text{ to } 100\text{ nF}$, $R_{\text{charge}} = 100\text{ k}\Omega$ ($I_{\text{charge}} = 50\,\mu\text{A}$).
    \item \textbf{Comparator \& Reset:} TL072 Schmitt trigger ($R_{\text{in}} = 10\text{ k}\Omega, R_{\text{fb}} = 60\text{ k}\Omega$), 2N3904 reset BJT.
    \item \textbf{LO Modulator:} Sweeps $f_{\text{LO}}$ from $20\text{ kHz}$ to $50\text{ kHz}$ synchronously with the $0\text{V} \to 5\text{V}$ tuning ramp.
\end{itemize}

\section{Comprehensive Simulation Test Suite}

\subsection{Sim 1: AC Frequency Response \& Stability Analysis}
\begin{lstlisting}
.ac dec 200 10k 50k
.meas AC f_center MAX mag(V(IF_out_2))
.meas AC f_low WHEN mag(V(IF_out_2))=f_center/sqrt(2) FALL=1
.meas AC f_high WHEN mag(V(IF_out_2))=f_center/sqrt(2) RISE=1
.meas AC bw PARAM (f_high - f_low)
\end{lstlisting}
\textbf{Target:} Verify $f_0 = 25,000 \pm 100\text{ Hz}$, $\text{BW} = 200 \pm 25\text{ Hz}$, and confirm phase decreasing monotonically with no right-half plane poles.

\subsection{Sim 2: Monte Carlo Component Tolerance Analysis}
\begin{lstlisting}
.param R3_val = {mc(796k, 0.01)}
.param R4_val = {mc(6.34k, 0.01)}
.param C3_val = {mc(1n, 0.05)}
.param C4_val = {mc(1n, 0.05)}
.step param run 1 100 1
\end{lstlisting}
\textbf{Target:} 100-run statistical sweep ensuring $>95\%$ of runs remain within $25,000 \pm 250\text{ Hz}$ ($1.0\%$) center frequency error.

\subsection{Sim 3: Worst-Case Sensitivity Analysis}
\begin{lstlisting}
.param R3_wc = {wc(796k, 0.01, 1)}
.param R5_wc = {wc(796k, 0.01, 2)}
.param C3_wc = {wc(1n, 0.05, 3)}
.param C4_wc = {wc(1n, 0.05, 4)}
.step param run 1 16 1
\end{lstlisting}
\textbf{Target:} Evaluate maximum frequency error bounding box ($\Delta f_{\text{max}}$) across all binary tolerance corner permutations.

\subsection{Sim 4: Thermal Drift Analysis}
\begin{lstlisting}
.model RES_PRECISION R (tc1=50e-6)  ; 50 ppm/C TCR
.step temp 0 70 10
\end{lstlisting}
\textbf{Target:} Ensure total frequency shift over $\Delta T = 40^\circ\text{C}$ is $< 50\text{ Hz}$ ($< 25\%$ of filter bandwidth).

\subsection{Sim 5: Mixer Intermodulation Distortion \& Linearity (FFT)}
\begin{lstlisting}
* Two-tone RF input: 5.0 kHz + 5.2 kHz (100 mVpk each); LO = 20.0 kHz
.tran 0 20m 10m 100n
\end{lstlisting}
\textbf{Target:} Take FFT on \texttt{V(mix\_out)}. Measure mixing products ($25.0, 25.2\text{ kHz}$) and 3rd-order intermodulation products ($24.8, 25.4\text{ kHz}$). Verify $\text{SFDR} \ge 45\text{ dBc}$.

\subsection{Sim 6: Master Dynamic Multi-Tone Sweep}
\begin{lstlisting}
* Three-tone RF: 4.0 kHz + 7.0 kHz + 12.0 kHz; Sweeping LO: 20 kHz to 40 kHz
.tran 0 100m 0 1u
\end{lstlisting}
\textbf{Target:} Plot $V_{\text{env\_out}}$ vs. $V(\text{ramp})$ in X-Y display mode. Confirm three distinct, symmetrical spectral peaks at exact expected positions.

\section{Summary of Deliverables}

\begin{table}[h!]
\centering
\small
\begin{tabularx}{\textwidth}{clllX}
\toprule
\textbf{\#} & \textbf{Schematic File} & \textbf{Simulation Type} & \textbf{Key Directive} & \textbf{Target Metric} \\
\midrule
1 & \texttt{03\_IF\_Filter\_Bank.asc} & AC Response & \texttt{.ac dec 200 10k 50k} & $f_0 = 25\text{ kHz}, \text{BW} = 200\text{ Hz}$, Stable \\
2 & \texttt{03\_IF\_Filter\_Bank.asc} & Monte Carlo & \texttt{.step param run 1 100 1} & $\Delta f_0 < \pm 250\text{ Hz}$ across 100 runs \\
3 & \texttt{03\_IF\_Filter\_Bank.asc} & Temperature & \texttt{.step temp 0 70 10} & $\Delta f_0 < 50\text{ Hz}$ over $\Delta T = 40^\circ\text{C}$ \\
4 & \texttt{02\_Gilbert\_Mixer.asc} & Linearity / FFT & \texttt{.tran 0 20m 10m 100n} & $\text{SFDR} \ge 45\text{ dBc}$, conversion gain \\
5 & \texttt{06\_Sawtooth\_VCO.asc} & Ramp Linearity & \texttt{.tran 0 50m 0 1u} & $0\text{V} \to 5\text{V}$ linear ramp, flyback $< 0.5\text{ ms}$ \\
6 & \texttt{SpectrumAnalyzer\_Master.asc} & Master Sweep & \texttt{.tran 0 100m 0 1u} & 3-tone spectrum resolved in X-Y mode \\
\bottomrule
\end{tabularx}
\caption{Required Simulation Files and Verification Milestones}
\end{table}

\end{document}
